\section{Weekly Summary}

\begin{tcolorbox}[colback=yellow!10,colframe=black!50,title=AI-Generated Content]
\noindent The summary below was written by Claude (Anthropic) from that week's regression outputs and series descriptions. It has not been reviewed or fact-checked by a human, and should be read as an automated, informal recap -- not analysis.
\end{tcolorbox}

Welcome back to the weekly ritual where we throw two unrelated economic series into a blender, hit "regress," and act surprised at whatever sludge comes out. Five days, five randomly selected odd couples, and -- spoiler alert -- the universe once again declined to hand us a Nobel Prize. As always, nothing here is a real finding. These are blind dates between spreadsheets. Any chemistry you see is almost certainly two numbers that both happened to go up during the same decade.

Monday kicked us off with Mid-Atlantic residential natural gas prices versus the percentage of syndicated loans at foreign bank branches, because of course those two belong in the same sentence. The raw regression served up a confident negative relationship with an R-squared of 0.44 and a satisfyingly tiny p-value -- genuinely the kind of number that would make an undergrad start writing a thesis. And then we detrended it, and the whole thing evaporated into a p-value of 0.56 and an R-squared of 0.02. Textbook trend mirage. Strip out the shared drift and these two go right back to being total strangers who nodded at each other once at a party. Tuesday did the exact same magic trick: bank equity ratios versus the dollar-to-SDR exchange rate looked "significant" raw, then collapsed to nothing once detrended. Two for two on the dunk tank.

Wednesday is where it got genuinely weird, in the fun way. Hong Kong's holdings of U.S. long-term corporate bonds versus commercial natural-gas carbon dioxide emissions -- a pairing so unhinged I want to frame it -- actually held up in both regressions. Raw R-squared around 0.24, detrended R-squared around 0.25, both significant, and the relationship even flipped sign after detrending for extra drama. This is the rare specimen that survives the trend-ectomy, which in this project counts as a sighting of Bigfoot. Do I think Hong Kong's bond appetite is secretly governed by American pipeline emissions? Absolutely not. But it's the most interesting coincidence we've bagged all week, so we're putting it on the mantel. Thursday, meanwhile, paired two flavors of discontinued bank-loan statistics and produced nothing significant in either version -- a merciful, boring result that I respect.

Friday closed things out with furniture-and-electronics retail inventories versus New England recreational-goods spending. The raw regression was significant, and -- plot twist -- it stayed significant after detrending too, with a p-value around 0.01. Slightly weaker than the raw version but still hanging on, which means "people's homes have stuff in them and people buy fun stuff" are at least loosely dancing to the same beat even after removing the shared march of time and inflation. Marginally plausible! Still not a claim I'd defend under oath. So this week's scoreboard: two spectacular trend mirages, one merciful nothingburger, and two relationships that stubbornly refused to die. Please remember that this entire exercise is random pairs of numbers being forced to hold hands, and spurious correlation is the house specialty, not the exception. See you next week for more statistically significant nonsense.
